% ---------- תקציר מנהלים ----------
% 22.11.2026: גוף שלישי, עד 200 מילים (הנחיית המרצה). כל מספר מגובה במקור בפרק 1.
\chapter*{תקציר}
\addcontentsline{toc}{chapter}{תקציר}

חברה ביטחונית-אווירית צריכה ספקים שמייצרים מוצרים איכותיים ואמינים, כי ספק חלש עלול לפגוע ברציפות האספקה של מערכות להגנת המדינה. בארגונים גדולים, קליטת ספק חדש נמשכת שבועות. הקושי גדול במיוחד בחברות פרטיות: שיעור הסיכון הפיננסי בהן גבוה יותר, ודווקא עליהן כמעט אין דוחות כספיים פומביים.

הדוח מציע מתודולוגיה להערכת ספק לפני ההתקשרות, שנשענת על מידע ציבורי בלבד ומבוצעת בידי סוכני בינה מלאכותית. ההערכה בנויה משני חלקים: \textbf{פרופיל חוזקות} (איכות ואמינות, התאמה טכנית, עמידה ביעדים ומצב פיננסי) ו\textbf{דגלי סיכון} (מצב גיאופוליטי, אירועים חריגים ואבטחת סייבר). ספק ברשימת סנקציות או ממדינה עוינת לישראל נפסל לפני כל בדיקה אחרת. המערכת מציגה גם כמה מהמידע הנדרש נמצא, וכשהוא חסר היא מפיקה רשימה לבדיקה ידנית. התוצר הוא המלצה מנומקת עם מקור לכל טענה, וההחלטה הסופית נשארת בידי אדם.

אבני הדרך: אב-טיפוס עובד בשבוע 10 של תוכנית העבודה, ותיקוף לאחור על מקרי כשל ציבוריים ועל ספקים מוכחים עד שבוע 13.

\vspace{0.5cm}
\noindent\textbf{מילות מפתח:} ניהול סיכוני ספקים; הערכת ספקים טרום התקשרות; מערכות מרובות סוכנים; מידע ממקורות גלויים (\foreignlanguage{english}{OSINT}); אבטחת סייבר בשרשרת האספקה; תעשייה ביטחונית-אווירית; מידע חסר.

\clearpage
